% Wrapper-free original proof. Requires the definitions and macros of
% bounded_transport.tex, input_output_incidence.tex and unary_boolean_transport.tex.
\section{Raw scales, prime incidence, and the original first-half hit code}

The source \emph{Transport into the colleague's stronger finite-code
positivity layers}, equations (PCT1)--(PCT11), was read in full.
The source \emph{Fixed TCHISLA rays and their exact intersection with
fixed divisor rays}, section \emph{The original factor scale and its
complete coordinate fibre}, was also read in full.  We retain their
raw scale and the PCT first-half radix.  The proofs here give their
exact composition with the bounded input--output incidence and with
the different full-range unary radix used earlier in this text.

\begin{theorem}[Raw factor scale with every original coordinate retained]
\label{thm:raw-scale-fibre}
Fix the original prime \(p=12h+1\), \(h\ge1\).
Let \(\mathcal T_p^{\rm raw}\) be the set of all tuples
\((A,B,C,D,\varepsilon)\) with \(A,B,C,D\in\Z_{>0}\) and
\(\varepsilon\in\{0,1\}\), satisfying
\[
 4ABCD=A+p^{1-\varepsilon}B+pC,
 \qquad a=ABD\in A_p=\{3h+1,\ldots,9h\}.
\]
Let \(\mathcal T_p^{\rm cop}\) be its sublocus \(\gcd(A,B)=1\).
There is a bijection with explicit inverse
\[
 \mathcal T_p^{\rm raw}\longleftrightarrow
 \{(A_0,B_0,C_0,D_0,\varepsilon,k):
       (A_0,B_0,C_0,D_0,\varepsilon)\in\mathcal T_p^{\rm cop},
       \ k>0,\ k^2\mid D_0\},
\]
given by
\[
 k=\gcd(A,B),\qquad
 (A_0,B_0,C_0,D_0)=(A/k,B/k,C/k,k^2D),
\]
and
\[
 (A,B,C,D)=(kA_0,kB_0,kC_0,D_0/k^2).
\]
The projection forgetting only \(k\) has the complete fibre
\(\{k>0:k^2\mid D_0\}\). Its size is
\[
 \prod_{\ell\mid D_0}(\lfloor v_\ell(D_0)/2\rfloor+1).
\]
The bijection retaining \(k\) has singleton fibres.

The following are equalities of the original ordered coordinates:
\begin{align*}
 a&=ABD=A_0B_0D_0,& R&=4a-p,&S&=pa,\\
 u&=B^2D=B_0^2D_0,&
 (x,y,z)&=(ABD,p^\varepsilon ACD,pBCD)\\
 &&&=(A_0B_0D_0,p^\varepsilon A_0C_0D_0,pB_0C_0D_0),\\
 d_z&=Rz-S=p^{2-\varepsilon}B^2D,
 &d_y&=Ry-S=p^\varepsilon A^2D,\\
 g_{\rm actual}&=\gcd(a,u)=kBD=B_0D_0.
\end{align*}
Thus every labelled cross residual \(4vw-p(v+w)\) for two
of the original denominators is unchanged.  The factor quotients
and affine numerators satisfy
\[
 \frac{B+p^\varepsilon C}{A}
   =\frac{B_0+p^\varepsilon C_0}{A_0},\quad
 BCD=B_0C_0D_0,\quad
 \frac{A+pC}{B}=\frac{A_0+pC_0}{B_0},
\]
\[
 \Omega=k\Omega_0,\quad\Theta=k\Theta_0,
 \quad \Omega+q\Theta=k(\Omega_0+q\Theta_0),
\]
where \(\Omega_0=A_0+\varepsilon B_0\) and
\(\Theta_0=C_0+(1-\varepsilon)B_0\).
The original ratio has reduced pair
\((A_0,p^{1-\varepsilon}B_0)\), whose residual gate quotient is
\(C_0=C/k\).  In particular none of these maps replaces a raw
tuple by an unlabelled coprime tuple.
\end{theorem}
\begin{proof}
Since \(A,B\le AB\le ABD=a<p\), one has \(1\le k<p\)
and \(\gcd(k,p)=1\).  The original equation modulo \(k\)
gives \(k\mid pC\); Bezout therefore gives \(k\mid C\).
All proposed coordinates are positive integers, and
\(\gcd(A_0,B_0)=1\).  The exact residual identity is
\[
 4ABCD-A-p^{1-\varepsilon}B-pC
 =k\bigl(4A_0B_0C_0D_0-A_0-p^{1-\varepsilon}B_0-pC_0\bigr).
\]
Its positive factor \(k\) shows equivalence of the two equations,
while \(ABD=A_0B_0D_0\) preserves the original shell bound.
Conversely the square-divisor condition makes \(D_0/k^2\) a
positive integer, the displayed residual vanishes, and
\(\gcd(kA_0,kB_0)=k\).  Both compositions therefore recover
every coordinate, including the retained scale.  Unique prime
factorization identifies its full fibre with
\(0\le v_\ell(k)\le\lfloor v_\ell(D_0)/2\rfloor\), proving
the count, including the empty prime product when \(D_0=1\).

Substitution proves the displayed equalities for the denominators,
\(a,R,S,u\), and all three factor quotients.  To check the oriented
divisor residuals, the original equation gives
\(RC=A+p^{1-\varepsilon}B\).  Multiplying this by \(pBD\)
and subtracting \(pABD\) gives \(Rz-S=p^{2-\varepsilon}B^2D\).
Multiplying instead by \(p^\varepsilon AD\) and subtracting
\(pABD\) gives \(Ry-S=p^\varepsilon A^2D\).
Also \(\gcd(ABD,B^2D)=BD\gcd(A,B)=kBD=B_0D_0\).
Each cross residual is unchanged because its two labelled
denominators and the original prime are unchanged.  The affine
numerator formulas follow by inserting \(kA_0,kB_0,kC_0\)
in their definitions.  Finally
\[
 \frac yz=\frac{A_0}{p^{1-\varepsilon}B_0},\qquad
 A_0+p^{1-\varepsilon}B_0=RC_0.
\]
The fraction is reduced because \(\gcd(A_0,B_0)=1\) and
\(A_0<p\), so \(p\nmid A_0\).  This proves the exact primitive
pair and the gate quotient without dropping its scale conversion.
In the original reconstruction \eqref{eq:reconstruct}, the coordinate
named \(d\) is precisely \(d_y=Ry-S\). The other oriented coordinate
is its explicit complement: \(d_y d_z=S^2\), as their displayed
products give \(p^2A^2B^2D^2=S^2\). Thus
\(d_z=S^2/d_y\), with the same inverse formula in the other direction;
no original orientation or name for \(d\) is reassigned.
The ray source's ordered labels \((d,e)\) correspond to
\((d_z,d_y)\); the original reconstruction order is
\((d_y,d_z)\). Their exact change of order is
\((r,s)\mapsto(s,r)\) on \(\{(r,s)\in\Z_{>0}^2:rs=S^2\}\),
whose square is the identity and whose fibres are singletons.
\end{proof}

\begin{theorem}[Exact scale restriction of the fixed primitive ray]
\label{thm:raw-scale-period}
Fix a coprime seed \((A_0,B_0,C_0,D_{0,p},\varepsilon)\)
in \(\mathcal T_p^{\rm cop}\).  Put
\[
 \begin{gathered}
 K_0=4A_0B_0C_0,\quad
 \Theta_0=C_0+(1-\varepsilon)B_0,\quad
 \Omega_0=A_0+\varepsilon B_0,\\
 g_0=\gcd(K_0,\Theta_0),\quad m_0=K_0/g_0,\quad w=\Theta_0/g_0.
 \end{gathered}
\]
Retain any \(k>0\) with \(k^2\mid D_{0,p}\), and the associated
fixed raw parameters \((kA_0,kB_0,kC_0,\varepsilon)\).
For every integer \(q\), keep the original rational functions
\[
 D_0(q)=\frac{\Omega_0+q\Theta_0}{K_0},\qquad
 D_k(q)=\frac{\Omega_0+q\Theta_0}{k^2K_0}.
\]
Their complete integrality statements are
\[
 D_0(q)\in\Z\ \Longleftrightarrow\ q\equiv p\pmod{m_0},
 \qquad
 D_k(q)\in\Z\ \Longleftrightarrow\ D_0(q)\in k^2\Z
 \Longleftrightarrow\ q\equiv p\pmod{m_k},
\]
where
\begin{equation}\label{eq:raw-scale-period}
 m_k=\frac{k^2K_0}{\gcd(k^2K_0,\Theta_0)},\qquad
 \frac{m_k}{m_0}=\frac{k^2}{\gcd(k^2,w)}.
\end{equation}
Thus the full raw return is the specified subray of the primitive
ray, not an enlargement of the raw domain.  Intersecting with
\(q\ge p\), \(q\equiv1\pmod{12}\) gives exactly
\[
 q=p+\lcm(12,m_k)t,\qquad t\in\Z_{\ge0}.
\]
Primality, and any further output predicates, remain exact tests on
these same integers.

For every such positive returned \(q\), the raw tuple and the
coprime tuple give the same original ordered integer witness and
the same \(a_q,R_q,S_q,u_q\).  The scale inverse is exactly
\(D_k(q)=D_0(q)/k^2\).  At a fixed primitive returned input \(q\),
the full set of scales coming from this fixed primitive seed is
\[
 \mathcal K_{p,q}=
 \{k>0:k^2\mid\gcd(D_{0,p},D_0(q))\}.
\]
It has the square-divisor count of this displayed gcd.
\end{theorem}
\begin{proof}
The primitive seed equation is
\(K_0D_{0,p}=\Omega_0+p\Theta_0\).  Therefore
\[
 D_0(q)=D_{0,p}+(q-p)\Theta_0/K_0.
\]
Because \(\gcd(m_0,w)=1\), its last summand is an integer
exactly when \(m_0\mid q-p\).  For the retained \(k\), the
seed value \(D_k(p)=D_{0,p}/k^2\) is already an integer, so
\[
 D_k(q)-D_k(p)=(q-p)\Theta_0/(k^2K_0)
\]
is integral exactly when \(m_k\mid q-p\), by division by the
positive gcd.  The equality \(D_0(q)=k^2D_k(q)\) proves both
directions of the square-divisibility statement, including when
the original primitive rational value is not an integer.

The period ratio follows by a complete gcd calculation:
\[
 \gcd(k^2K_0,\Theta_0)
 =g_0\gcd(k^2m_0,w)=g_0\gcd(k^2,w).
\]
The last equality follows prime by prime, since no prime divides
both \(m_0\) and \(w\).  Substitution gives
\eqref{eq:raw-scale-period}.  Its ratio is a positive integer
dividing \(k^2\), so every raw return also passes the primitive
period.  Since the seed satisfies \(p\equiv1\pmod{12}\),
the simultaneous congruences are exactly divisibility of \(q-p\)
by \(\lcm(12,m_k)\); \(q\ge p\) gives the stated parameter
and its inverse \((q-p)/\lcm(12,m_k)\).

For \(q\ge p>0\), the numerator \(\Omega_0+q\Theta_0\) is
positive, hence each integral returned \(D_k(q)\) is positive.
Apply the coordinate identities of Theorem~\ref{thm:raw-scale-fibre}
at this input: their direct substitutions are valid for the fixed
positive returned integer \(q\) and preserve all listed coordinates.
The original identity at \(q\) gives
\[
 \frac1{A_0B_0D_0(q)}+
 \frac1{q^\varepsilon A_0C_0D_0(q)}+
 \frac1{qB_0C_0D_0(q)}=\frac4q,
\]
by the same common denominator computation as
\eqref{eq:incidence-witness}.  If \(q\) is prime and
\(q\equiv1\pmod{12}\), its first denominator lies in the original
full \(A_q\): the exact formula
\[
 \frac{a_q}{q}=\frac{\Theta_0}{4C_0}
                       +\frac{\Omega_0}{4C_0q}
\]
is greater than \(1/4\) and at most \(a_p/p\), which is at most
\(3/4-3/(4p)\le3/4-3/(4q)\).
Integrality then gives the original full integer endpoints.
Finally a retained scale must satisfy its seed condition
\(k^2\mid D_{0,p}\) and its exact return condition
\(k^2\mid D_0(q)\); these are equivalent to their displayed gcd
divisibility.  Prime factorization gives its complete count and
the scale reconstruction inverse.
\end{proof}

\begin{theorem}[First-half gate and exact change of unary radix]
\label{thm:raw-first-half-code}
Use the fixed primitive seed and ray above, and a returned prime
\(q\ge p\), \(q\equiv1\pmod{12}\).
Put \(a_q=A_0B_0D_0(q)\), \(R_q=4a_q-q\), \(S_q=qa_q\).
The exact first-half return condition is
\begin{equation}\label{eq:raw-first-half-gate}
 2a_q\le q-1
 \quad\Longleftrightarrow\quad
 q(2C_0-\Theta_0)\ge\Omega_0+2C_0.
\end{equation}
If \(2C_0-\Theta_0\le0\), no positive input passes it.
Otherwise it is precisely
\[
 q\ge\left\lceil\frac{\Omega_0+2C_0}{2C_0-\Theta_0}\right\rceil.
\]
This gate is invariant under every retained raw scale.  If the
original seed is already first-half, it passes for every returned
\(q\ge p\); otherwise the computed gate remains in the domain.

On this gate retain the PCT radix, with names distinguished from
both the input CRT modulus and the full-range radix:
\[
 M_H=(q-1)/2,\quad J_H=(q-1)/4,\quad
 \Lambda_H=(q-1)^2/4=2M_HJ_H,\quad N_H=M_H\Lambda_H.
\]
The exact PCT hit of this ray is
\begin{equation}\label{eq:raw-pct-code}
 j=(4a_q-q-3)/4,\qquad
 c_H=(A_0-1)\Lambda_H+(q-1)j+2(B_0-1)+\varepsilon.
\end{equation}
Here \(0\le j<J_H\), \(1\le A_0,B_0\le M_H\), and
\[
 A_0B_0\mid a_q,\quad \gcd(A_0,B_0)=1,\quad
 R_q\mid A_0+q^{1-\varepsilon}B_0,
\]
with exact quotients \(D_0(q)\) and \(C_0\).

For every code in this PCT rectangle, its inverse is the following
successive division, with all remainders retained:
\[
 A=1+\lfloor c_H/\Lambda_H\rfloor,\quad
 v=c_H-(A-1)\Lambda_H,\quad j=\lfloor v/(q-1)\rfloor,
\]
\[
 r=v-(q-1)j,\qquad B=1+\lfloor r/2\rfloor,
 \qquad\varepsilon=r-2\lfloor r/2\rfloor.
\]
For every \(0\le c<N_H\), let
\((A(c),j(c),B(c),\varepsilon(c))\) be that exact decoder and put
\[
 R(c)=4j(c)+3,\qquad a(c)=(q+R(c))/4.
\]
Define the total bit \(\chi_H(c)\) to be one precisely when
\[
 A(c)B(c)\mid a(c),\qquad \gcd(A(c),B(c))=1,\qquad
 R(c)\mid A(c)+q^{1-\varepsilon(c)}B(c),
\]
and zero otherwise. Define its total least-hit selector by the original product
\[
 c_{*,H}(q)=\sum_{n=0}^{N_H-1}\prod_{c=0}^n(1-\chi_H(c)).
\]
Then the displayed witness gives
\[
 0\le c_{*,H}(q)\le c_H<A_0\Lambda_H.
\]
For every integer \(1\le r_0\le M_H\), the existence of a hit
with \(A\le r_0\) is equivalent to
\(c_{*,H}(q)<r_0\Lambda_H\).
The value \(N_H\) occurs exactly when there is no hit and is
never decoded as a candidate.  A specified witness has
\(A_0\le2\) exactly when its own code is below
\(2\Lambda_H\); this implies that same least-hit bound but its
failure does not exclude another hit in the first two layers.

Write \(q=12h_q+1\) and retain separately the earlier full radix
\[
 M_F=9h_q,\quad J_F=6h_q,\quad\Lambda_F=2M_FJ_F,
 \quad N_F=M_F\Lambda_F.
\]
For an arbitrary full-range code \(c_F\), decode
\((j,A,B,\varepsilon)\) in the full radix and define its own
shell \(a(c_F)=(q+4j+3)/4\).
There is an exact bijection from full-range hit codes satisfying
\(2a(c_F)\le q-1\) onto PCT hit codes, given by
\[
 c_F\longmapsto(A-1)\Lambda_H+(q-1)j+2(B-1)+\varepsilon.
\]
Its inverse decodes the PCT radix and
uses \((A-1)\Lambda_F+2M_Fj+2(B-1)+\varepsilon\).
Both maps have singleton fibres and preserve the full tuple,
\(a_q,R_q,S_q,C,D,u\), the original ordered witness, and every
predicate on those retained coordinates.  A raw scale is attached
separately; forgetting it has exactly the scale fibre already proved.
Their exact numerical difference is
\[
 \Lambda_F=3\Lambda_H,\qquad
 c_F-c_H=2(A-1)\Lambda_H+\frac{q-1}{2}j.
\]
\end{theorem}
\begin{proof}
The exact equality \(a_q=(\Omega_0+q\Theta_0)/(4C_0)\)
makes \(2a_q\le q-1\) equivalent, after multiplying by the
positive integer \(2C_0\), to
\(\Omega_0+q\Theta_0\le2C_0q-2C_0\).
This is \eqref{eq:raw-first-half-gate}.  Its right side requires
a positive number \(\Omega_0+2C_0\); hence a nonpositive
coefficient \(2C_0-\Theta_0\) permits no positive \(q\).
For a positive coefficient, integer division gives exactly the
displayed ceiling.  If the seed passes at \(p\), that coefficient
is positive and the inequality remains true for every \(q\ge p\).
The first denominator is unchanged by raw scaling, proving gate
invariance without an extra range assumption.

The full-shell bounds proved above give \(a_q>q/4\), so the
integer \(R_q=4a_q-q\equiv3\pmod4\) is at least three.
The first-half bound gives \(R_q\le q-2\), and consequently
\[
 0\le j=(R_q-3)/4\le(q-5)/4=J_H-1.
\]
Also \(a_q=A_0B_0D_0(q)\le M_H\), so
\(A_0,B_0\le M_H\) and \(A_0B_0\mid a_q\).
The primitive tuple equation at \(q\) gives
\(R_qC_0=A_0+q^{1-\varepsilon}B_0\), with \(C_0>0\).
These prove precisely the three hit tests and their quotients.

For arbitrary rectangle data the digit
\(2(B-1)+\varepsilon\) runs through
\(0,\ldots,2M_H-1=q-2\) exactly once.
Together with \(j=0,\ldots,J_H-1\) it gives every remainder
\(0,\ldots,(q-1)J_H-1=\Lambda_H-1\) exactly once.
The outer digit \(A-1\) therefore codes exactly
\(0,\ldots,N_H-1\), and the displayed successive divisions
recover all four original digits.  In particular the ray's hit
has code less than \(A_0\Lambda_H\).
If the least hit is \(c_0\), the product defining the selector
is one for precisely \(n<c_0\) and zero for \(n\ge c_0\).
Its sum is \(c_0\); if none exists, all \(N_H\) terms are one.
Thus the proved hit gives the claimed upper bound.  The first
\(r_0\) layers consist exactly of codes below \(r_0\Lambda_H\).
Their occupancy is therefore equivalent to the least-hit inequality,
including the failure sentinel case.  The same digit bound proves
the statement about the specified first-two-layer witness.

For the change of radix, a full-range hit with first-half shell has
\(j<J_H=3h_q\) and \(A,B\le a_q\le M_H=6h_q\)
by its divisor condition.  Hence its decoded digits belong to the
PCT rectangle.  Conversely all PCT hit digits lie in the full
rectangle because \(J_H\le J_F\), \(M_H\le M_F\), and their
first denominator is \(3h_q+j+1\le6h_q=M_H\).
The shell and hit tests are identical on these same digits.
Both decoding algorithms were proved inverse to their respective
radix formulas, so decode--recode in either order is the identity
on the digits and on the corresponding codes.  On a hit the exact
quotients are \(D=a_q/(AB)\),
\(C=(A+q^{1-\varepsilon}B)/R_q\), and \(u=B^2D\).
These are computed from the same preserved digits and shell.
The ordered denominator formulas consequently agree in both
presentations.  A predicate on those coordinates has the same truth
value under the bijection.  The raw scale is not one of the code
digits; its complete inverse choices remain those of
Theorems~\ref{thm:raw-scale-fibre} and \ref{thm:raw-scale-period}.
Finally \(M_F=9h_q\), \(J_F=6h_q\), \(M_H=6h_q\),
\(J_H=3h_q\) give \(\Lambda_F=108h_q^2=3\Lambda_H\).
Subtracting the two complete code formulas retains the same inner
digit and gives
\(c_F-c_H=2(A-1)\Lambda_H+(18h_q-12h_q)j\),
which is exactly the displayed difference.
\end{proof}

\begin{theorem}[The complete scaled incidence of first-half hits]
\label{thm:raw-scale-hit-incidence}
Fix the primitive seed of Theorem~\ref{thm:raw-scale-period},
the exact bounded input set \(\mathcal X\) of
Theorem~\ref{thm:input-output-incidence}, and its original input
data, bounds and output modulus \(\mathcal D\).
Put \(\mathcal K_p=\{k>0:k^2\mid D_{0,p}\}\).
The complete scaled first-half incidence is
\[
 \mathcal I_H^{\rm raw}=
 \left\{(k,X,q):
 \begin{array}{l}
 k\in\mathcal K_p,\quad X\in\mathcal X,\quad q\text{ prime},
       \quad p<q\le N(V),\quad q\mid N(X),\\
 q\equiv1\pmod{12},\quad q\equiv1\pmod{\mathcal D},
       \quad q\equiv p\pmod{m_k},\\
 q(2C_0-\Theta_0)\ge\Omega_0+2C_0
 \end{array}\right\}.
\]
It has a bijection onto the graph recording, in addition to all
fixed original data and \(k,X,q\), the following complete output:
\[
 \begin{gathered}
 D_0(q),\quad
 (A,B,C,D,\varepsilon)
       =(kA_0,kB_0,kC_0,D_0(q)/k^2,\varepsilon),\\
 (a_q,y_q,z_q)=(A_0B_0D_0(q),
      q^\varepsilon A_0C_0D_0(q),qB_0C_0D_0(q)),\\
 R_q=4a_q-q,\quad S_q=qa_q,\quad u_q=B_0^2D_0(q),
 \quad j=(4a_q-q-3)/4,\quad c_H,\quad c_F.
 \end{gathered}
\]
Here \(c_H,c_F\) are the two distinct codes of
Theorem~\ref{thm:raw-first-half-code}.  The graph inverse is
projection to the retained fixed data and \(k,X,q\), and its
fibres are singletons.

For a fixed \(k\), the exact projection to \((X,q)\) is the
input--output incidence theorem with the raw period \(m_k\),
restricted by the computed first-half gate.  In particular for
every retained \(q\) its input fibre is the disjoint root-branch
union \eqref{eq:incidence-root-branch}, with precisely the same
original \(a_{\rm in},M,U,V\); every original input constraint
remains imposed.  For fixed \((k,X)\), its fibre is exactly the
distinct prime divisors of \(N(X)\) satisfying all the displayed
output tests for that \(k\).

Let \(\mathcal I_H^0\) be the primitive incidence, obtained by
using \(m_0\) and the same first-half gate and other constraints.
Forgetting \(k\) maps \(\mathcal I_H^{\rm raw}\) onto
\(\mathcal I_H^0\), with exact fibre
\[
 \mathcal K_{p,q}
 =\{k>0:k^2\mid\gcd(D_{0,p},D_0(q))\}
 \quad\text{over }(X,q).
\]
This projection has a specified section \(k=1\); it does not
assert that one fixed raw \(k>1\) returns on the entire primitive
ray.  Consequently the two exact bounded counts are
\[
 |\mathcal I_H^{\rm raw}|
 =\sum_{(X,q)\in\mathcal I_H^0}
   \prod_{\ell\mid\gcd(D_{0,p},D_0(q))}
       \left(\left\lfloor
          \frac{v_\ell(\gcd(D_{0,p},D_0(q)))}2
        \right\rfloor+1\right),
\]
and the sum over \(k\in\mathcal K_p\), retained prime \(q\),
and compatible roots \(\rho\) of the branch sizes in
\eqref{eq:incidence-count}, using that \(k\)'s output tests.

Every retained output proves
\(c_{*,H}(q)\le c_H<A_0\Lambda_H\).
Any finite collection of exact decidable arithmetic predicates
and Boolean expressions may be evaluated on this one retained
graph tuple.  Its truth fibres correspond bijectively under the
graph and radix inverses.  Under a projection forgetting
\(k\) or \(X\), the full selected fibre is exactly the fibre
above intersected with those predicate tests; no predicate is
replaced by independent witnesses for its different occurrences.
The union, intersection, fibre-product and count formulas of
Theorem~\ref{thm:unary-truth-fibres} therefore apply on this same
finite graph domain.  No assertion of prime coverage or universal
first-layer occupancy is required for any of these correspondences.
\end{theorem}
\begin{proof}
The scale set is finite by its proved square-divisor description.
Each retained scale gives a valid original raw seed with its
original first denominator unchanged.  Its exact output period is
\(m_k\), by Theorem~\ref{thm:raw-scale-period}.
Applying the input--output incidence theorem to these fixed raw
parameters uses
\[
 \frac{4(kA_0)(kB_0)(kC_0)}
      {\gcd(4(kA_0)(kB_0)(kC_0),k\Theta_0)}
 =\frac{k^2K_0}{\gcd(k^2K_0,\Theta_0)}=m_k.
\]
Thus this is exactly its raw integrality condition, including every
scale factor.  Its root congruence concerns only the retained input
polynomial \(N(X)\), so the root and CRT input branches retain
the same original input data and all their inverse parameters.
The first-half restriction is precisely the inequality already
proved and is a predicate on that same \(q\).
These observations prove the fixed-scale fibres in both projection
directions, including every empty fibre.

Every such output has \(D_0(q)/k^2\) a positive integer,
and the coordinate equalities of the scale theorem give exactly
the listed raw tuple, original ordered witness and shell data.
The gate gives the PCT code and its full-radix counterpart with
their proved inverses and least-hit inequality.
All graph coordinates are uniquely determined from the retained
fixed seed and \(k,X,q\); projection recovers these labels.
This proves both inverse compositions and the singleton graph
fibres without identifying inputs that happen to yield the same
prime or witness.

Forgetting \(k\) preserves a primitive incidence because
\(m_0\mid m_k\) and all other tests are unchanged.
Conversely a primitive incidence has \(D_0(q)\) a positive
integer.  The scales lifting it are exactly those satisfying both
\(k^2\mid D_{0,p}\) and \(k^2\mid D_0(q)\), by the proved
raw period equivalence.  This is the displayed gcd fibre.
It is nonempty because \(k=1\) always lies in it, proving
surjectivity and the stated section.  Unique prime factorization
counts that fibre.  Counting all triples by this projection gives
the first count; counting first by \(k\), then \(q\), then the
disjoint input root branches gives the second.  All sets are finite
by the finite scale set, bounded prime set and finite lift intervals.

Finally a specified predicate is evaluated at the exact graph
coordinates that have just been proved integral and recovered by
the inverse.  Reindexing either code retains those coordinates, so
it retains the predicate's truth value.  A projection with a
nontrivial fibre retains precisely those members that pass the same
predicate tests; its selected fibre is therefore the stated set
intersection.  Every graph element has one complete truth vector.
The disjoint truth-fibre decomposition, diagonal intersection
fibre product, union multiplicities and finite inclusion--exclusion
proof of Theorem~\ref{thm:unary-truth-fibres} applies to these
finite sets verbatim, with the graph element as the common variable.
This establishes all Boolean correspondences and counts without
assuming that any prescribed fibre is occupied.
\end{proof}

\paragraph{Exact radix and retained-scale fixtures.}
At the prime \(p=13\), the primitive tuple
\((A_0,B_0,C_0,D_{0,p},\varepsilon)=(2,1,5,2,0)\)
satisfies \(4\cdot2\cdot1\cdot5\cdot2=80=2+13+13\cdot5\).
Its original first denominator is \(a=4\), residual \(R=3\)
and index \(j=0\).  The PCT radix has
\((M_H,J_H,\Lambda_H)=(6,3,36)\), so \(c_H=36\).
The full radix has \((M_F,J_F,\Lambda_F)=(9,6,108)\),
so the exact reindexed code is \(c_F=108\).
If the unchanged integer \(36\) were used as a full code, its
full decoder would instead give \((A,j,B,\varepsilon)=(1,2,1,0)\),
with \(a=6\), \(R=11\).  It fails because \(11\nmid1+13\).
This checks the numerical distinction together with its exact
tuple-preserving reindexing map.

For a separate scale fixture, retain \(p=61\),
\((A,B,C,D,\varepsilon)=(2,4,2,2,1)\).
The equation is \(128=2+4+61\cdot2\) and \(a=16\)
lies in \(A_{61}=\{16,\ldots,45\}\) and in its first half.
Its scale is \(k=2\), and the primitive tuple is \((1,2,1,8,1)\).
Thus \(K_0=8\), \(\Theta_0=1\), \(\Omega_0=3\),
\(m_0=8\), \(m_2=32\), with the respective modulo-twelve
periods \(24\) and \(96\).
The prime \(q=109=61+2\cdot24\) gives
\[
 D_0(109)=14,\qquad a_{109}=28,\qquad R_{109}=3,
 \qquad j=0,\qquad c_H=2(2-1)+1=3.
\]
Its ordered primitive witness is \((28,1526,3052)\).
Over denominator \(3052\), its reciprocal numerators are
\(109,2,1\), whose sum \(112\) gives \(4/109\).
The first-half hit is therefore actual, but the retained raw
coefficient is \(D_2(109)=14/4=7/2\), not an integer.
The scale fibre over this returned prime has only \(k=1\),
since \(\gcd(D_{0,p},D_0(109))=\gcd(8,14)=2\).
The prime \(q=157=61+96\) instead gives
\(D_0(157)=20\), \(D_2(157)=5\), and the actual raw tuple
\((2,4,2,5,1)\), with the same ordered witness
\((40,3140,6280)\) as its primitive tuple \((1,2,1,20,1)\).
Here the scale fibre is \(\{1,2\}\), since \(\gcd(8,20)=4\).
Primality of \(61,109,157\) follows by trial division by the
primes up to their square roots: respectively
\(2,3,5,7\), \(2,3,5,7\), and \(2,3,5,7,11\);
none divides its indicated number.  These computations exhibit the
exact strict restriction of a fixed raw scale without enlarging its
ray when a primitive code succeeds.
